\section{El sistema de numeración}

Imagine que quiere contar piedras en una bolsa, pero todavía no dispone de nombres ni de símbolos para expresar cuántas hay. Una forma sencilla de hacerlo sería asociar un dedo a cada piedra. Este recurso permite llevar un registro mientras la cantidad de piedras no supere el número de dedos disponibles. La idea importante es que, para contar y comunicar cantidades, necesitamos alguna forma de representarlas.

Un sistema de numeración es un conjunto de símbolos y reglas que permite representar números. Los números no deben confundirse con los símbolos que usamos para escribirlos: una misma cantidad puede tener representaciones diferentes según el sistema elegido. La elección de los símbolos y de la forma de agruparlos es una convención desarrollada históricamente, y no existe una única manera posible de hacerlo.

En la actualidad se utilizan distintos sistemas de numeración. Algunos de ellos son:
\begin{itemize}
  \item Sistema binario: utiliza dos dígitos, \(0\) y \(1\).
  \item Sistema octal: utiliza ocho dígitos, \(0,1,2,3,4,5,6,7\).
  \item Sistema decimal: utiliza diez dígitos, \(0,1,2,3,4,5,6,7,8,9\). Es el sistema de numeración de uso cotidiano más extendido.
  \item Sistema hexadecimal: utiliza dieciséis dígitos, \(0,1,2,3,4,5,6,7,\) \(8,9,A,B,C,D,E,F\), y es muy empleado en computación.
  \item Sistema sexagesimal: organiza las cantidades en grupos de sesenta. En la medición del tiempo todavía aparece esta forma de agrupamiento: por ejemplo, \(60\) segundos forman \(1\) minuto y \(60\) minutos forman \(1\) hora.
\end{itemize}

Existen muchos otros sistemas de numeración. Lo esencial es distinguir la cantidad que se quiere expresar de la representación elegida para escribirla. Al contar objetos, puede pensarse intuitivamente que se establece una correspondencia entre cada objeto y una marca, una palabra o un numeral; el sistema de numeración proporciona una manera compartida y organizada de realizar esa representación.

\subsection{El sistema decimal}

Este es el sistema que más va a utilizar, el sistema decimal. Este sistema consta de diez dígitos cada uno con un nombre:
\begin{itemize}
  \item \(0\): ``cero'' \(\rightarrow\) representa cantidad nula. Para el ejemplo de contar piedras significa que no hay piedras.
  \item \(1\): ``uno'' \(\rightarrow\) representa una unidad. Este número es importante, porque a partir de él y una operación llamada suma se obtiene el resto de cantidades.
  \item \(2\): ``dos'' \(\rightarrow\) representa dos unidades. Aunque esto pueda parecerle absurdo y trivial, matemáticamente hablando es una verdad. Para que usted (lector) no quede con un vacío existencial después de leer sobre números, puede pensarlo de la siguiente manera: si definimos un dedo (por ejemplo el índice) como unidad, luego podemos decir que tenemos \(2\) dedos índice. Uno en cada mano, porque cada uno representa una unidad, y dos unidades se representarán con el símbolo 2, o con dos dedos.
  \item \(3\): ``tres'' \(\rightarrow\) representa tres unidades. Análogo a lo que sucede con el número 2, el tres es el dos más una unidad. Y esto continúa así infinitamente.
\end{itemize}
Este sistema se llama decimal porque tiene en total diez dígitos individuales. Estos dígitos pueden combinarse para representar cantidades mayores a \(9\). Aquí, por ejemplo, el número \(17\) es una representación con dos números de una cantidad mayor a \(9\). Cuando llega a \(9\), en el sistema decimal, se le han agotado los símbolos disponibles, de esta manera puede intuir que esto es una gran limitación si pretende contar cualquier cosa. Para salvar esta problemática, se usa la combinación de varios dígitos, asignando la cantidad más significativa al número de la izquierda y la menor al número de la derecha.\footnote{Combinatoria se refiere a mezclar números donde su orden de posición es importante, es decir, 12 no es lo mismo que 21, ya que la cifra más significativa en el primer caso es 1 y en el segundo caso 2. Esto significa que el segundo número es mayor.} Utilizar más números agrupados para representar cantidades tan grandes como se necesite es la solución.

Los números a partir del los 90s tienen nombres muy bien definidos, pasando por las centenas (cien, doscientos, \dots), luego millares (mil, diez mil, cien mil, \dots) y así cada número tendrá un nombre asociado. 

De este modo, si se dice ciento cincuenta y ocho mil cuatrocientos treinta y tres, se sabe que se refiere específicamente al número 158433. Los nombres de los números pasando los millones, billones y trillones tienen nombres muy peculiares. Por ejemplo, el \href{https://es.wikipedia.org/wiki/G\%C3\%BAgol\#G.C3.BAgolplex}{gugol} (ver enlace\footnote{\href{https://es.wikipedia.org/wiki/G\%C3\%BAgol\#G.C3.BAgolplex}{\texttt{https://es.wikipedia.org/wiki/G\%C3\%BAgol\#G.C3.BAgolplex}}}). Estos números tan grandes no son muy usuales de nombrar, pero son divertidos de conocer.
